\scanpage{14}\mspage{13}
% ink top, opening and lemma
Donc il faut simplement vérifier le

\uline{Lemme} Soit $\mathcal L$ un ens ordonné, et $x,y,u,v$ trois \unsure{élémts|éléments}, tels que
% ink upper third, clauses a–d and brace at right
\[
\begin{aligned}
&\text{a)}\quad u,v\text{ sont entre }x\text{ et }y\\
&\left.\begin{aligned}
&\text{b)}\quad u\text{ est entre }x\text{ et }v\\
&([\,\text{c)}\quad v\text{ est entre }u\text{ et }y\,]\deleted{\unsure{)}}
\end{aligned}\ \right\}\quad
\begin{aligned}&\text{une des deux}\\&\text{en suffit}\end{aligned}\\
&\text{d)}\quad x<y\quad,
\end{aligned}
\]
alors on a
\[
x\leq u\leq v\leq y
\]

% ink middle, Dém, first four lines
\uline{Dém} \deleted{\unsure{L'ass|L'assert}} \unsure{les|Les} conditions b) c) d) impliquent que $I\deleted{\unsure{\{x}}=\{x,y,u,v\}$ est un drapeau, \deleted{\unsure{Prouvons|Prenons}} et \unsure{SPG} $\mathcal L=I$ est un ens tot. ordonné \deleted{que $x$ en} à 2, 3, ou 4 \unsure{élémts|éléments}.
% ink middle, crossed proof and interlinear replacement
\deleted{Alors \unsure{Prouvons|Prenons}} \inserted{Comme $u,v$ sont} \deleted{que $x$ est le plus petit élément, donc (comme $x<y$) que l'on ne peut avoir $u<x$, ni $v<x$. Supposons $u<x$ ($<y$) ; comme $u$ est entre $x$ et $v$, on aurait \unsure{$v\leq u<$}\ $v<u<x<y$} \inserted{entre $x$ et $y$ et $x<y$,} on doit avoir $x\leq u,v\leq y$, \deleted{il reste à prouver $u\leq v$} et comme $u$ est entre $x$ et $v$, on doit avoir $x\leq u\leq v$\quad OK

% ink bottom, NB and lettered lines
\uline{NB} On a utilisé
\[
\begin{aligned}
&\text{a) la divisibilité de la rel. d'ordre}\\
&\text{b) le fait que }\mathcal L\text{ soit filtrant}\\
&\hspace{2em}\text{croissant et filtrant décroissant.}
\end{aligned}
\]
% ink last two lines
Et a) n'a servi qu'à montrer que la relation ``$a$ et $b$ disjoints'' est déductible de ``$a$ est str. entre $b$ et $c$''

% pale vertical watermark
\rotatebox{90}{\textsc{\unsure{g Extra Strong}}}
